% GENERATED FILE. Edit canonical lessons, then run npm run generate:books.
% canonical-source-hash: fnv1a64:1970a35ec2fff48c
% canonical-lessons: MW-C239-juvar, MW-C239-sarsu, MW-C239-kapas, MW-C239-badhno-grow, MW-C239-kupal

\chapter{Crops}
\label{ch:mw-crops}
\hlchaptermodality{239}

\begin{chapteropening}
\textbf{By the end of this chapter:} \emph{I can name the crops of the desert farm and say they grow.}

Complete the last lesson of chapter 239: I can name the crops of the desert farm and say they grow.
\end{chapteropening}

\section[\texorpdfstring{\mw{जुवार} (juvār)}{जुवार (juvār)}]{\mw{जुवार} (juvār) — sorghum; the great millet}

\label{lesson:MW-C239-juvar}



\begin{quote}
Before the new one: say the Marwadi for a ploughman, then the Marwadi for a seed.
\end{quote}

\subsection*{You'll want to know: \mw{जुवार}}
\textbf{\mw{जुवार}} — \emph{juvār} — \textquotedblleft{}sorghum; the great millet\textquotedblright{}.

\begin{quote}
\mw{खेत} \mw{मांय} \mw{जुवार} \mw{हरी} \mw{है।}
\end{quote}

\begin{quote}
\emph{khet mā̃y juvār harī hai.}
\end{quote}

\begin{quote}
The sorghum in the field is green.
\end{quote}

A feminine noun.

\begin{cousinweb}
\textbf{\mw{जुवार}} (\emph{juvār}) is inherited from Sanskrit \emph{yavākāra}, \textquotedblleft{}shaped like barley\textquotedblright{}; CDIAL lists Old Marwari \emph{juvāri}, \textquotedblleft{}millet\textquotedblright{} (CDIAL 10437). In Marwadi it names sorghum, the tall millet, beside \textbf{\mw{बाजरो}}.
\end{cousinweb}

\begin{grammarlens}[title={What you've built}]
One more word for this chapter.
\end{grammarlens}

\subsection*{Your turn}
\begin{itemize}
\raggedright
  \item \emph{Say it:} \emph{juvār}
  \item \emph{Say it:} \emph{juvār}, once more
  \item \emph{Say it:} \emph{khet mā̃y juvār harī hai.} — The sorghum in the field is green.
  \item \emph{Recall:} say the Marwadi for a ploughman, then the Marwadi for a seed, then say \emph{juvār} again
\end{itemize}

\begin{tcolorbox}[breakable,colback=teal!4,colframe=teal!35!black,title={Before you move on}]
What does \mw{जुवार} mean? (\textquotedblleft{}Sorghum; the great millet\textquotedblright{}.) Say it once more.
\end{tcolorbox}
\section[\texorpdfstring{\mw{सरसूं} (sarsū̃)}{सरसूं (sarsū̃)}]{\mw{सरसूं} (sarsū̃) — mustard (the plant and its seed)}

\label{lesson:MW-C239-sarsu}



\begin{quote}
Before the new one: say the Marwadi for a seed, then the Marwadi for sorghum.
\end{quote}

\subsection*{You'll want to know: \mw{सरसूं}}
\textbf{\mw{सरसूं}} — \emph{sarsū̃} — \textquotedblleft{}mustard (the plant and its seed)\textquotedblright{}.

\begin{quote}
\mw{सरसूं} \mw{पीळी} \mw{है।}
\end{quote}

\begin{quote}
\emph{sarsū̃ pīḷī hai.}
\end{quote}

\begin{quote}
The mustard is yellow.
\end{quote}

A feminine noun.

\begin{cousinweb}
\textbf{\mw{सरसूं}} (\emph{sarsū̃}) is inherited from Sanskrit \emph{sarṣapa}, \textquotedblleft{}mustard seed\textquotedblright{}; CDIAL lists Hindi \emph{sarsõ}, the mustard plant. CDIAL gives no Marwari form in this entry (CDIAL 13281).
\end{cousinweb}

\begin{grammarlens}[title={What you've built}]
One more word for this chapter.
\end{grammarlens}

\subsection*{Your turn}
\begin{itemize}
\raggedright
  \item \emph{Say it:} \emph{sarsū̃}
  \item \emph{Say it:} \emph{sarsū̃}, once more
  \item \emph{Say it:} \emph{sarsū̃ pīḷī hai.} — The mustard is yellow.
  \item \emph{Recall:} say the Marwadi for a seed, then the Marwadi for sorghum, then say \emph{sarsū̃} again
\end{itemize}

\begin{tcolorbox}[breakable,colback=teal!4,colframe=teal!35!black,title={Before you move on}]
What does \mw{सरसूं} mean? (\textquotedblleft{}Mustard (the plant and its seed)\textquotedblright{}.) Say it once more.
\end{tcolorbox}
\section[\texorpdfstring{\mw{कपास} (kapās)}{कपास (kapās)}]{\mw{कपास} (kapās) — cotton (the plant)}

\label{lesson:MW-C239-kapas}



\begin{quote}
Before the new one: say the Marwadi for sorghum, then the Marwadi for mustard.
\end{quote}

\subsection*{You'll want to know: \mw{कपास}}
\textbf{\mw{कपास}} — \emph{kapās} — \textquotedblleft{}cotton (the plant)\textquotedblright{}.

\begin{quote}
\mw{हूं} \mw{कपास} \mw{बेचूं} \mw{हूं।}
\end{quote}

\begin{quote}
\emph{hū̃ kapās bechū̃ hū̃.}
\end{quote}

\begin{quote}
I sell cotton.
\end{quote}

A feminine noun.

\begin{cousinweb}
\textbf{\mw{कपास}} (\emph{kapās}) is inherited from Sanskrit \emph{karpāsa}, \textquotedblleft{}the cotton plant\textquotedblright{}, which CDIAL takes, perhaps, from Austroasiatic; CDIAL lists Hindi and Gujarati \emph{kapās}. CDIAL gives no Marwari form in this entry (CDIAL 2877).
\end{cousinweb}

\begin{grammarlens}[title={What you've built}]
One more word for this chapter.
\end{grammarlens}

\subsection*{Your turn}
\begin{itemize}
\raggedright
  \item \emph{Say it:} \emph{kapās}
  \item \emph{Say it:} \emph{kapās}, once more
  \item \emph{Say it:} \emph{hū̃ kapās bechū̃ hū̃.} — I sell cotton.
  \item \emph{Recall:} say the Marwadi for sorghum, then the Marwadi for mustard, then say \emph{kapās} again
\end{itemize}

\begin{tcolorbox}[breakable,colback=teal!4,colframe=teal!35!black,title={Before you move on}]
What does \mw{कपास} mean? (\textquotedblleft{}Cotton (the plant)\textquotedblright{}.) Say it once more.
\end{tcolorbox}
\section[\texorpdfstring{\mw{बधणो} (badhṇo)}{बधणो (badhṇo)}]{\mw{बधणो} (badhṇo) — to grow; to increase}

\label{lesson:MW-C239-badhno-grow}



\begin{quote}
Before the new one: say the Marwadi for mustard, then the Marwadi for cotton.
\end{quote}

\subsection*{You'll want to know: \mw{बधणो}}
\textbf{\mw{बधणो}} — \emph{badhṇo} — \textquotedblleft{}to grow; to increase\textquotedblright{}.

\begin{quote}
\mw{बरसात} \mw{मांय} \mw{बाजरो} \mw{बधै} \mw{है।}
\end{quote}

\begin{quote}
\emph{barsāt mā̃y bājro badhai hai.}
\end{quote}

\begin{quote}
The millet grows in the rains.
\end{quote}

Past (it …): \mw{वो} \mw{रूंख} \mw{बध्यो}, \emph{vo rū̃kh badhyo}.

\begin{cousinweb}
\textbf{\mw{बधणो}} (\emph{badhṇo}) is inherited from Sanskrit \emph{vardhate}, \textquotedblleft{}grows\textquotedblright{}; CDIAL lists the Old Marwari verb \emph{badhaï}, \textquotedblleft{}to increase, to thrive\textquotedblright{} (CDIAL 11376).
\end{cousinweb}

\begin{grammarlens}[title={What you've built}]
One more word for this chapter.
\end{grammarlens}

\subsection*{Your turn}
\begin{itemize}
\raggedright
  \item \emph{Say it:} \emph{badhṇo}
  \item \emph{Say it:} \emph{badhṇo}, once more
  \item \emph{Say it:} \emph{barsāt mā̃y bājro badhai hai.} — The millet grows in the rains.
  \item \emph{Recall:} say the Marwadi for mustard, then the Marwadi for cotton, then say \emph{badhṇo} again
\end{itemize}

\begin{tcolorbox}[breakable,colback=teal!4,colframe=teal!35!black,title={Before you move on}]
What does \mw{बधणो} mean? (\textquotedblleft{}To grow; to increase\textquotedblright{}.) Say it once more.
\end{tcolorbox}
\section[\texorpdfstring{\mw{कूंपळ} (kū̃paḷ)}{कूंपळ (kū̃paḷ)}]{\mw{कूंपळ} (kū̃paḷ) — a new shoot; a bud}

\label{lesson:MW-C239-kupal}



\begin{quote}
Before the new one: say the Marwadi for cotton, then the Marwadi for to grow.
\end{quote}

\subsection*{You'll want to know: \mw{कूंपळ}}
\textbf{\mw{कूंपळ}} — \emph{kū̃paḷ} — \textquotedblleft{}a new shoot; a bud\textquotedblright{}.

\begin{quote}
\mw{रूंख} \mw{माथै} \mw{नुवी} \mw{कूंपळ} \mw{है।}
\end{quote}

\begin{quote}
\emph{rū̃kh māthai nuvī kū̃paḷ hai.}
\end{quote}

\begin{quote}
There's a new shoot on the tree.
\end{quote}

A feminine noun.

\begin{cousinweb}
\textbf{\mw{कूंपळ}} (\emph{kū̃paḷ}) is inherited from Sanskrit \emph{kuḍmala}, \textquotedblleft{}a bud\textquotedblright{}; CDIAL lists Old Marwari \emph{kūpaḷa}, \textquotedblleft{}a fresh bud or shoot\textquotedblright{} (CDIAL 3250).
\end{cousinweb}

\begin{grammarlens}[title={What you've built}]
One more word for this chapter.
\end{grammarlens}

\subsection*{Your turn}
\begin{itemize}
\raggedright
  \item \emph{Say it:} \emph{kū̃paḷ}
  \item \emph{Say it:} \emph{kū̃paḷ}, once more
  \item \emph{Say it:} \emph{rū̃kh māthai nuvī kū̃paḷ hai.} — There's a new shoot on the tree.
  \item \emph{Recall:} say the Marwadi for cotton, then the Marwadi for to grow, then say \emph{kū̃paḷ} again
\end{itemize}

\begin{tcolorbox}[breakable,colback=teal!4,colframe=teal!35!black,title={Before you move on}]
What does \mw{कूंपळ} mean? (\textquotedblleft{}A new shoot; a bud\textquotedblright{}.) Say it once more.
\end{tcolorbox}
